\documentclass[12pt,a4paper]{article}

\usepackage[T2A]{fontenc}
\usepackage[utf8]{inputenc}
\usepackage[russian]{babel}
\usepackage{amsmath}
\usepackage{amssymb}
\usepackage{booktabs}
\usepackage{geometry}
\usepackage{pgfplots}

\geometry{margin=2cm}
\pgfplotsset{compat=1.18}

\title{Физическая модель падения воды на минивен}
\author{Микропроект по вычислительной физике}
\date{}

\begin{document}

\maketitle

\begin{abstract}
В работе рассматривается учебная численная модель удара воды массой
\(2\) тонны по упрощённой модели кузова автомобиля. Вода моделируется методом
SPH, бокс строится и размешивается в Gmsh, а его деформация рассчитывается в
FEniCSx. Модель не претендует на точный инженерный расчёт реального минивэна,
но позволяет качественно воспроизвести падение воды, передачу импульса, смятие
тонкостенной конструкции и визуализацию процесса в ParaView.
\end{abstract}

\section{Постановка задачи}

Идея проекта состоит в том, чтобы смоделировать падение большого объёма воды на
кузов автомобиля. Полная геометрия минивэна на этом этапе заменена более простой
моделью: тонкостенным полым параллелепипедом из стали. Такое упрощение полезно,
потому что оно оставляет главную физику удара, но делает задачу понятной и
управляемой.

Вода имеет массу
\[
M_w = 2000 \ \text{кг}.
\]
Если принять плотность воды равной
\[
\rho_0 = 1000 \ \text{кг}/\text{м}^3,
\]
то соответствующий объём воды равен
\[
V_w = \frac{M_w}{\rho_0}
=
\frac{2000}{1000}
=
2 \ \text{м}^3.
\]

В коде этот объём задаётся прямоугольным блоком:
\[
x \in [-1,1], \quad
y \in [-0.5,0.5], \quad
z \in [50,51].
\]

\section{Оценка падения без сопротивления воздуха}

Перед численным моделированием полезно оценить порядок величин. Если сначала
забыть про сопротивление воздуха, то скорость падения можно найти из закона
сохранения энергии:
\[
M_w g H = \frac{M_w v^2}{2}.
\]
Отсюда
\[
v = \sqrt{2gH}.
\]

При \(g=9.81 \ \text{м}/\text{с}^2\) и \(H=50 \ \text{м}\):
\[
v = \sqrt{2 \cdot 9.81 \cdot 50}
\approx 31.3 \ \text{м}/\text{с}.
\]

Это примерно
\[
31.3 \ \text{м}/\text{с}
\approx 113 \ \text{км}/\text{ч}.
\]

Время свободного падения с высоты \(50\) метров:
\[
t = \sqrt{\frac{2H}{g}}
=
\sqrt{\frac{100}{9.81}}
\approx 3.19 \ \text{с}.
\]

В самой модели нижняя граница воды находится на высоте \(z=50\), а крыша бокса
находится примерно на высоте
\[
z_{\text{roof}} = 1.8 \ \text{м}.
\]
Поэтому расстояние от нижней части воды до крыши:
\[
H_{\text{roof}} = 50 - 1.8 = 48.2 \ \text{м}.
\]
Оценка времени до первого контакта:
\[
t_{\text{roof}} =
\sqrt{\frac{2H_{\text{roof}}}{g}}
=
\sqrt{\frac{96.4}{9.81}}
\approx 3.13 \ \text{с}.
\]

Это хорошо согласуется с численным расчётом: первые заметные нагрузки на крышу
появляются около \(t \approx 3.2\) секунды.

\begin{figure}[h]
\centering
\begin{tikzpicture}
\begin{axis}[
    width=0.82\textwidth,
    height=6cm,
    grid=both,
    xlabel={\(t\), с},
    ylabel={\(z\), м},
    xmin=0, xmax=3.4,
    ymin=0, ymax=52,
    legend pos=north east
]
\addplot[blue, thick, domain=0:3.2, samples=100] {50 - 0.5*9.81*x^2};
\addlegendentry{низ блока воды}
\addplot[red, dashed, thick, domain=0:3.4] {1.8};
\addlegendentry{исходная крыша бокса}
\end{axis}
\end{tikzpicture}
\caption{Оценка положения нижней части воды при свободном падении. Пересечение
с уровнем крыши происходит примерно за \(3.13\) секунды.}
\end{figure}

\section{Энергия и импульс удара}

Потенциальная энергия воды на высоте \(50\) метров:
\[
E_p = M_w g H.
\]
Подставим числа:
\[
E_p = 2000 \cdot 9.81 \cdot 50
= 981000 \ \text{Дж}
\approx 0.98 \ \text{МДж}.
\]

Если почти вся эта энергия переходит в кинетическую энергию перед ударом, то
удар получается очень сильным. Это объясняет, почему в реальном эксперименте
кузов автомобиля буквально разнесло.

Импульс воды перед ударом можно оценить так:
\[
P_w = M_w v.
\]
При \(v \approx 31.3 \ \text{м}/\text{с}\):
\[
P_w \approx 2000 \cdot 31.3
= 62600 \ \text{кг}\cdot\text{м}/\text{с}.
\]

В программе именно импульс оказался важнее одного только пика силы. Если
ориентироваться только на мгновенный пик силы, то первые отдельные частицы могут
слишком рано включить смятие бокса. Поэтому в модели учитывается накопленный
импульс за кадр.

\begin{table}[h]
\centering
\begin{tabular}{lll}
\toprule
Величина & Формула & Значение \\
\midrule
Масса воды & \(M_w\) & \(2000 \ \text{кг}\) \\
Объём воды & \(M_w/\rho_0\) & \(2 \ \text{м}^3\) \\
Время падения с 50 м & \(\sqrt{2H/g}\) & \(3.19 \ \text{с}\) \\
Скорость перед ударом & \(\sqrt{2gH}\) & \(31.3 \ \text{м}/\text{с}\) \\
Энергия & \(M_wgH\) & \(0.98 \ \text{МДж}\) \\
Импульс & \(M_wv\) & \(6.26\cdot 10^4 \ \text{кг}\cdot\text{м}/\text{с}\) \\
\bottomrule
\end{tabular}
\caption{Оценки основных величин до удара.}
\end{table}

\section{Геометрия и масса бокса}

Бокс имеет внешние размеры:
\[
L_x = 4.0 \ \text{м}, \quad
L_y = 2.0 \ \text{м}, \quad
L_z = 1.8 \ \text{м}.
\]

Толщина стенки:
\[
t = 0.0008 \ \text{м}.
\]

Плотность стали:
\[
\rho_s = 7800 \ \text{кг}/\text{м}^3.
\]

Внешний объём параллелепипеда:
\[
V_{\text{out}} = L_xL_yL_z
= 4.0 \cdot 2.0 \cdot 1.8
= 14.4 \ \text{м}^3.
\]

Внутренние размеры после вырезания полости:
\[
L_x^{in}=L_x-2t=3.9984 \ \text{м},
\]
\[
L_y^{in}=L_y-2t=1.9984 \ \text{м},
\]
\[
L_z^{in}=L_z-2t=1.7984 \ \text{м}.
\]

Внутренний объём:
\[
V_{\text{in}} \approx 14.36994 \ \text{м}^3.
\]

Объём стали:
\[
V_s = V_{\text{out}} - V_{\text{in}}
\approx 14.4 - 14.36994
\approx 0.03006 \ \text{м}^3.
\]

Масса стальной оболочки:
\[
M_s = \rho_s V_s
\approx 7800 \cdot 0.03006
\approx 234.5 \ \text{кг}.
\]

Это число важно для понимания масштаба: на оболочку массой порядка
\(235\) кг падает вода массой \(2000\) кг. Масса воды примерно в
\[
\frac{2000}{234.5} \approx 8.5
\]
раз больше массы тонкостенного бокса.

\section{Дискретизация воды}

Вода разбивается на частицы (метод SPH). В текущей реализации расстояние между начальными
частицами:
\[
\Delta x = 0.15 \ \text{м}.
\]

Так как блок воды имеет размеры \(2 \times 1 \times 1\) м, при таком шаге
получается:
\[
N_x = 14, \quad N_y = 7, \quad N_z = 7.
\]

Общее число частиц:
\[
N = N_xN_yN_z
= 14 \cdot 7 \cdot 7
= 686.
\]

Масса одной частицы:
\[
m = \frac{M_w}{N}
=
\frac{2000}{686}
\approx 2.92 \ \text{кг}.
\]

Это не означает, что в реальности вода состоит из шариков по \(2.92\) кг.
Это численная аппроксимация: каждая частица представляет небольшой объём воды.

\begin{table}[h]
\centering
\begin{tabular}{lll}
\toprule
Параметр & Обозначение & Значение в коде \\
\midrule
Шаг по времени & \(\Delta t\) & \(0.004 \ \text{с}\) \\
Число шагов & \(N_t\) & \(2000\) \\
Время расчёта & \(T=N_t\Delta t\) & \(8 \ \text{с}\) \\
Частота сохранения & -- & каждый 20-й шаг \\
Шаг частиц & \(\Delta x\) & \(0.15 \ \text{м}\) \\
Радиус сглаживания SPH & \(h\) & \(0.28 \ \text{м}\) \\
Число частиц & \(N\) & \(686\) \\
Масса частицы & \(m\) & \(2.92 \ \text{кг}\) \\
\bottomrule
\end{tabular}
\caption{Численные параметры SPH-части модели.}
\end{table}

\section{SPH-модель воды}

SPH расшифровывается как Smoothed Particle Hydrodynamics, то есть метод
сглаженных частиц. Главная идея: жидкость не хранится на фиксированной сетке.
Вместо этого она задаётся набором частиц, которые движутся вместе с потоком.

Для каждой частицы \(i\) хранятся:
\[
r_i, \quad v_i, \quad a_i, \quad \rho_i, \quad p_i.
\]
Здесь \(r_i\) -- положение, \(v_i\) -- скорость, \(a_i\) -- ускорение,
\(\rho_i\) -- плотность, \(p_i\) -- давление.

\subsection{Плотность}

Плотность частицы считается суммой вкладов соседних частиц:
\[
\rho_i = \sum_j m_j W(|r_i-r_j|, h).
\]

Функция \(W\) называется ядром сглаживания. Она устроена так, что близкие
частицы влияют сильно, а далёкие частицы не влияют совсем. В коде используется
Poly6-ядро:
\[
W(r,h)=
\begin{cases}
\dfrac{315}{64\pi h^9}(h^2-r^2)^3, & 0 \le r < h,\\
0, & r \ge h.
\end{cases}
\]

Если расстояние между частицами больше \(h=0.28\) м, то вклад соседа равен
нулю. Это делает расчёт локальным.

\subsection{Давление}

Давление считается по простой линейной формуле:
\[
p_i = k(\rho_i-\rho_0).
\]

В коде используется
\[
k = 120.
\]
Это не настоящий модуль сжимаемости воды. Для реальной воды он был бы намного
больше. Здесь \(k\) выбран маленьким, чтобы учебная симуляция оставалась
устойчивой и не требовала очень маленького шага по времени.

В реализации используется ограничение:
\[
p_i = \max(0, k(\rho_i-\rho_0)).
\]

Причина простая: если разрешить большое отрицательное давление, частицы на
свободной поверхности могут начать неестественно притягиваться друг к другу.
Для текущей визуальной модели это только ухудшает результат.

\subsection{Сила давления}

Сила давления для частицы \(i\):
\[
F_i^{p} =
- \sum_j m_j
\frac{p_i+p_j}{2\rho_j}
\nabla W_{\text{spiky}}(|r_i-r_j|,h).
\]

Знак минус означает, что частицы выталкиваются из областей повышенного давления.
Если частицы сильно сблизились, плотность растёт, давление растёт, и сила
старается раздвинуть их.

Градиент Spiky-ядра:
\[
\nabla W_{\text{spiky}}(r,h)
=
-\frac{45}{\pi h^6}(h-r)^2 \frac{r_i-r_j}{|r_i-r_j|}.
\]

\subsection{Вязкость}

Вязкость отвечает за сглаживание разницы скоростей:
\[
F_i^{\nu} =
\mu \sum_j m_j
\frac{v_j-v_i}{\rho_j}
\nabla^2 W_{\nu}(|r_i-r_j|,h).
\]

В текущей реализации:
\[
\mu = 24.
\]

Это тоже эффективный численный коэффициент, а не точная динамическая вязкость
воды. Он нужен, чтобы поток не распадался на слишком шумные отдельные частицы.

\subsection{Сопротивление воздуха}

Сопротивление воздуха задано простой линейной моделью:
\[
F_i^{air} = -c_{air} v_i.
\]

В коде:
\[
c_{air}=0.05.
\]

Такая модель не описывает воздух как отдельную среду. Она только слегка гасит
скорости частиц и делает движение спокойнее.

\subsection{Итоговое уравнение движения частиц}

Полное ускорение частицы:
\[
a_i =
g
+
\frac{F_i^p + F_i^\nu + F_i^{air}}{\rho_i}.
\]

После этого используется явная схема:
\[
v_i^{n+1} = v_i^n + a_i^n \Delta t,
\]
\[
r_i^{n+1} = r_i^n + v_i^{n+1}\Delta t.
\]

Схема простая, зато её легко читать и проверять.

\section{Контакт воды с боксом}

Бокс для SPH-части задаётся как набор внешних плоскостей. Когда частица попадает
внутрь объёма бокса, она выталкивается на ближайшую грань, а её скорость
меняется.

Нормальная скорость после удара:
\[
v_n^{new} = -e v_n^{old}.
\]

Для крыши выбран маленький коэффициент восстановления:
\[
e_{\text{roof}} = 0.03.
\]
Это означает почти неупругий удар: вода должна в основном терять вертикальную
скорость, а не отскакивать вверх.

Для боковых стенок:
\[
e_{\text{wall}} = 0.10.
\]

Касательная скорость тоже гасится:
\[
v_{\tau}^{new} = (1-d_{\tau})v_{\tau}^{old}.
\]

Для крыши:
\[
d_{\tau}^{roof}=0.55,
\]
то есть касательная скорость уменьшается на \(55\%\).

Импульс, переданный одной частицей:
\[
J_i = m_i |v_n^{new}-v_n^{old}|.
\]

Если частица массой \(2.92\) кг ударяется о крышу со скоростью около
\(30\ \text{м}/\text{с}\), то грубая оценка импульса одной частицы:
\[
J_i \sim 2.92 \cdot 30
\approx 87.6 \ \text{кг}\cdot\text{м}/\text{с}.
\]

Если таких частиц много, суммарный импульс быстро становится большим.
Поэтому в CSV-файл записывается не только максимальная сила, но и импульс за
интервал кадра:
\[
J_{\text{roof}} = \sum_i J_i.
\]

Эквивалентная средняя сила за интервал:
\[
F_{\text{avg}} \approx \frac{J_{\text{roof}}}{\Delta t_{\text{frame}}}.
\]

Такой подход устойчивее, чем использование только одного пика силы.

\section{Почему вода не должна зависать на крыше}

Есть важная особенность реализации. SPH-вода считается в C++, а деформация бокса
через FEniCSx считается после этого как постобработка. Поэтому вода напрямую не
видит деформированную FEniCSx-сетку.

Если ничего не исправлять, возникает неприятный визуальный эффект: бокс уже
смят в ParaView, но частицы воды продолжают останавливаться на старой плоскости
крыши.

Для этого в SPH-контакт добавлена простая модель опускания контактной крыши:
\[
s_{\text{roof}}^{new}
=
\min(s_{\max}, s_{\text{roof}}^{old} + \alpha J_{\text{roof}}).
\]

В текущей реализации:
\[
\alpha = 2.0\cdot 10^{-5},
\]
\[
s_{\max}=1.45 \ \text{м}.
\]

Это значит, что после сильного удара контактная плоскость крыши постепенно
уходит вниз. Следующие частицы воды сталкиваются уже не с исходной крышей
\(z=1.8\), а с опущенной плоскостью. Поэтому вода может попадать в область
смятого кузова, и в ParaView она больше не выглядит зависшей в воздухе.

\begin{figure}[h]
\centering
\begin{tikzpicture}
\begin{axis}[
    width=0.82\textwidth,
    height=6cm,
    grid=both,
    xlabel={номер кадра},
    ylabel={характерная высота воды, м},
    xmin=39, xmax=60,
    ymin=0, ymax=2.1,
    legend pos=north east
]
\addplot[blue, mark=*] coordinates {
    (39,2.195)
    (40,0.707)
    (41,0.628)
    (45,0.000)
    (50,0.000)
    (60,0.000)
};
\addlegendentry{\(z_{\min}\) воды}
\addplot[red, dashed, thick, domain=39:60] {1.8};
\addlegendentry{исходная крыша}
\end{axis}
\end{tikzpicture}
\caption{После удара частицы воды уходят ниже исходной высоты крыши. Это
показывает, что вода больше не фиксируется на старой, ещё не смятой плоскости.}
\end{figure}

\section{Контакт воды с землёй}

Земля задаётся плоскостью:
\[
z=0.
\]

При контакте с землёй вертикальная скорость воды почти полностью гасится:
\[
v_z^{new} \approx 0, \quad \text{если } v_z < 0.
\]

Горизонтальная скорость уменьшается из-за трения:
\[
v_x^{new} = (1-f)v_x,
\]
\[
v_y^{new} = (1-f)v_y.
\]

В коде:
\[
f = 0.18.
\]

Это означает, что после каждого контакта с землёй горизонтальная скорость
уменьшается примерно на \(18\%\). Такая земля не является настоящей моделью
грунта, но она не даёт воде бесконечно проваливаться вниз и уменьшает
нереалистичный отскок.

\section{Деформация бокса в FEniCSx}

Сетка бокса строится в Gmsh и передаётся в FEniCSx. Там решается задача линейной
упругости:
\[
-\nabla \cdot \sigma(u)=f.
\]

Здесь:
\[
u = (u_x,u_y,u_z)
\]
-- поле перемещений узлов сетки, а \(f\) -- нагрузка от удара воды.

Тензор малых деформаций:
\[
\varepsilon(u)=
\frac{1}{2}\left(\nabla u + \nabla u^T\right).
\]

Физический смысл: \(\varepsilon(u)\) показывает, насколько материал растягивается,
сжимается или сдвигается около каждой точки.

Для изотропного материала используется закон Гука:
\[
\sigma(u)=2\mu\varepsilon(u)+\lambda \operatorname{tr}(\varepsilon(u))I.
\]

Коэффициенты Ламе:
\[
\mu=\frac{E}{2(1+\nu)},
\]
\[
\lambda=\frac{E\nu}{(1+\nu)(1-2\nu)}.
\]

Для настоящей стали модуль Юнга порядка \(2\cdot 10^{11}\) Па. В проекте
используется пониженный эффективный модуль:
\[
E_{\text{eff}} = 1.5\cdot 10^9 \ \text{Па}.
\]

Это сделано осознанно. Текущая модель не содержит настоящей пластичности,
потери устойчивости и разрушения тонких стенок. Если поставить реальный модуль
стали без этих эффектов, бокс будет слишком жёстким и почти не будет сминаться.
Пониженный \(E_{\text{eff}}\) грубо имитирует потерю устойчивости тонкой
оболочки после удара.

Коэффициент Пуассона:
\[
\nu = 0.30.
\]

Тогда:
\[
\mu =
\frac{1.5\cdot 10^9}{2(1+0.3)}
\approx 5.77\cdot 10^8 \ \text{Па},
\]
\[
\lambda =
\frac{1.5\cdot 10^9\cdot 0.3}{(1+0.3)(1-2\cdot 0.3)}
\approx 8.65\cdot 10^8 \ \text{Па}.
\]

\section{Слабая форма задачи}

В FEniCSx решается не дифференциальное уравнение напрямую, а его слабая форма.
Она записывается так:
\[
\int_{\Omega}\sigma(u):\varepsilon(v)\,dx
=
\int_{\Gamma_{\text{roof}}} t_{\text{roof}}\cdot v\,ds
+
\int_{\Gamma_{\text{side}}} t_{\text{side}}\cdot v\,ds.
\]

Левая часть отвечает за внутреннее сопротивление материала деформации.
Правая часть отвечает за внешние нагрузки от воды.

Нижняя грань бокса закрепляется:
\[
u=0 \quad \text{на нижней грани}.
\]

Это приближённая модель контакта с землёй: днище не может свободно улететь вниз,
потому что оно опирается на плоскость.

\section{Остаточное смятие}

Обычная линейная упругость возвращает тело обратно после снятия нагрузки. Для
удара двух тонн воды это выглядит плохо: кузов должен остаться смятым. Поэтому
после упругого решения добавляется простая модель остаточной деформации.

На каждом кадре измеряются максимальные прогибы:
\[
d_{\text{roof}}, \quad d_{\text{side}}, \quad d_{\text{floor}}.
\]

Дальше хранятся накопленные значения:
\[
D_{\text{roof}}^{new}
=
\max(D_{\text{roof}}^{old}, 0.92d_{\text{roof}}),
\]
\[
D_{\text{side}}^{new}
=
\max(D_{\text{side}}^{old}, 0.88d_{\text{side}}),
\]
\[
D_{\text{floor}}^{new}
=
\max(D_{\text{floor}}^{old}, 0.55d_{\text{floor}}).
\]

Числа \(0.92\), \(0.88\), \(0.55\) задают, какая часть максимального прогиба
остаётся как необратимое смятие. Крыша сохраняет деформацию сильнее всего,
потому что основной удар приходит сверху.

Итоговая форма бокса:
\[
x_{\text{new}}
=
x_{\text{mesh}}
+
u_{\text{elastic}}
+
u_{\text{plastic}}.
\]

Здесь:

\begin{itemize}
  \item \(x_{\text{mesh}}\) -- исходные узлы сетки Gmsh;
  \item \(u_{\text{elastic}}\) -- упругое перемещение из FEniCSx;
  \item \(u_{\text{plastic}}\) -- добавленное остаточное смятие.
\end{itemize}

\begin{figure}[h]
\centering
\begin{tikzpicture}
\begin{axis}[
    width=0.82\textwidth,
    height=6cm,
    grid=both,
    xlabel={номер кадра},
    ylabel={деформация, м},
    xmin=39, xmax=100,
    ymin=0, ymax=2.0,
    legend pos=north east
]
\addplot[blue, mark=*] coordinates {
    (39,0.000)
    (40,1.830)
    (45,1.711)
    (60,1.706)
    (100,1.706)
};
\addlegendentry{крыша}
\addplot[orange, mark=square*] coordinates {
    (39,0.000)
    (40,1.063)
    (45,0.953)
    (60,0.952)
    (100,0.952)
};
\addlegendentry{стенки}
\end{axis}
\end{tikzpicture}
\caption{Пример поведения остаточной деформации в текущем расчёте. После удара
бокс не возвращается к исходной форме, а остаётся смятым.}
\end{figure}

\section{Как всё связано в программе}

Расчёт в программе идёт по этапам.

\begin{enumerate}
  \item В C++ задаются параметры воды, бокса и земли.
  \item Gmsh строит тонкостенную тетраэдральную сетку бокса.
  \item SPH-часть считает движение частиц воды.
  \item При контакте воды с боксом считаются силы и импульсы удара.
  \item Эти нагрузки записываются в файл \texttt{box\_loads.csv}.
  \item Python-скрипт с FEniCSx читает сетку и нагрузки.
  \item FEniCSx решает задачу упругости для каждого кадра.
  \item К упругому решению добавляется остаточное смятие.
  \item Вода, земля и деформированный бокс записываются в VTK-форматы.
  \item Результат просматривается в ParaView.
\end{enumerate}

\section{Что модель учитывает}

В текущей версии учитываются:

\begin{itemize}
  \item гравитационное падение воды;
  \item движение воды методом SPH;
  \item давление и вязкость воды;
  \item простое сопротивление воздуха;
  \item контакт воды с боксом;
  \item контакт воды с землёй;
  \item тетраэдральная сетка бокса из Gmsh;
  \item упругая деформация бокса в FEniCSx;
  \item остаточное смятие бокса после удара;
  \item визуализация в ParaView.
\end{itemize}

\section{Ограничения модели}

Модель остаётся учебной. В ней пока не учитываются:

\begin{itemize}
  \item настоящая нелинейная пластичность стали;
  \item разрушение, трещины и отрыв элементов;
  \item самоконтакт смятых стенок;
  \item воздух как отдельная расчётная среда;
  \item сложная модель грунта;
  \item точная геометрия реального минивэна.
\end{itemize}

Главное назначение этой версии -- получить понятную качественную модель. Она
показывает, как можно связать SPH для воды, Gmsh для сетки, FEniCSx для
деформации и ParaView для визуализации.

\end{document}
